\section*{ID7813: Linear Stability and Bifurcation (How Cymatics Create Vortices): Ψ(x,t)}
\paragraph{Metadata.}
PiID: 2150784245167 | RTEID: RTE:P30339.6141:T6.0081 | UUID: e1fd267e-c0f5-532b-b97e-3dc7ac242522

\paragraph{Canonical Statement.}
Assume base state \textbackslash{}(\textbackslash{}Psi(\textbackslash{}mathbf\{x\},t)=\textbackslash{}Psi\_0 + \textbackslash{}delta\textbackslash{}psi(\textbackslash{}mathbf\{x\},t)\textbackslash{}), with \textbackslash{}(\textbackslash{}Psi\_0\textbackslash{}) constant satisfying: \textbackslash{}[ 0 = -\textbackslash{}alpha\textbackslash{}Psi\_0 + g\textbackslash{}bar p\textbackslash{},\textbackslash{}Psi\_0 - \textbackslash{}beta|\textbackslash{}Psi\_0|\textasciicircum{}2\textbackslash{}Psi\_0, \textbackslash{}] where \textbackslash{}(\textbackslash{}bar p\textbackslash{}) is time-average pressure. Solve for \textbackslash{}(|\textbackslash{}Psi\_0|\textasciicircum{}2\textbackslash{}).

\paragraph{Canonical Equation.}
\[
\Psi(\mathbf{x},t)=\Psi_0 + \delta\psi(\mathbf{x},t)
\]

\paragraph{Dictionary.}
Relation: Ψ(x,t)=Ψ\_0 + δψ(x,t)

\paragraph{Encyclopedia.}
Linear Stability and Bifurcation (How Cymatics Create Vortices): Ψ(x,t) is an algebraic definition, written Ψ(x,t)=Ψ\_0 + δψ(x,t). It is an algebraic definition expressing the quantity directly in terms of the variables on its right-hand side. It is built from δ, ψ, x, defining Ψ(x,t). Within the Trawin Zero Point registry it serves as one element of the framework's mathematical narrative.
